\documentclass[11pt]{ctexart}
\usepackage{amsmath,amssymb,amsthm}
\usepackage{booktabs}
\usepackage[margin=1in]{geometry}
\usepackage{hyperref}
\hypersetup{pdftitle={M1 非停机问题的归约推导},pdfauthor={Z.ai}}

\newtheorem{theorem}{定理}[section]
\newtheorem{lemma}[theorem]{引理}
\newtheorem{proposition}[theorem]{命题}
\newtheorem{corollary}[theorem]{推论}
\theoremstyle{definition}
\newtheorem{definition}[theorem]{定义}
\newtheorem{assumption}[theorem]{计算假设}
\newtheorem{remark}[theorem]{注}
\newtheorem{conjecture}[theorem]{猜想}

\newcommand{\vtwo}[1]{\nu_2(#1)}
\newcommand{\HALT}{\textsf{HALT}}

\title{\bfseries M1（\texttt{1RB1LA\_0RC0RD\_1LC1LA\_1RE1RD\_1LF0LA---0LE}）的停机性\\ \large 从相位演算到记录时刻表的完全归约推导}
\author{}
\date{2026 年 10 月}

\begin{document}
\maketitle

\begin{abstract}
本文给出 6 态图灵机 $M=$ \texttt{1RB1LA\_0RC0RD\_1LC1LA\_1RE1RD\_1LF0LA---0LE}（空白带启动）非停机问题的完整归约推导：从相位演算出发，经锚点族的序曲定理、地标族轮定律的强归纳架构，把``$M$ 永不停机''归约为单一计算假设（主控封闭，假设 M），再经签名动力学的有限状态分析，把该假设归约为自主 2-adic 游走记录时刻表的延续，并证明其 32-位切片经有限证书无环且终止（全深度非最终周期性为猜想，K$>$32）。另证两个无条件于假设 M 的定理：全轨道模 3 台账（$\Sigma\equiv0\pmod3$，自第 54 步起；Conway 型迫使链入口 $S(3,1,3)$ 与一切 $\Sigma\not\equiv0$ 构型整半空间不可达，重入/锚点同余定理化），与停机等价（$M$ 停机 $\iff$ 某次 S-偶出生值 $A=x+y/2+1$ 为 2 的幂，双向代数证明；轨道侧 $10^7$ 事件 2,908 次出生零 2-幂为 B 级佐证）；出生守卫的逐深度污染有至深度 3 的裸机见证（见证级，非定理）。核心代数推导（扫掠、减半复合、序曲、类参数化、阶梯良基性、装配）完整给出；支配引理（D）的余项上界与粗模板不可能性（必要精度）分别降为 A$^-$ 与猜想（证据与已核部分见正文注）；机器验证部分以声明域穷尽的形式显式标出（台账 A/B/C 三级）。
\end{abstract}

\section{引言}

设 $M$ 为转移表
\[
\begin{array}{c|cc}
 & \text{读 }0 & \text{读 }1\\\hline
A & 1RB & 1LA\\
B & 0RC & 0RD\\
C & 1LC & 1LA\\
D & 1RE & 1RD\\
E & 1LF & 0LA\\
F & \HALT & 0LE
\end{array}
\]
（记号 \texttt{sXY}：写 $s$、移动 $X$、转入状态 $Y$；$F$ 读 $0$ 未定义为唯一停机转移。）本文证明的归约链概括为
\[
\text{$M$ 永不停机}\;\Longleftarrow\;\text{假设 M}\;\Longleftrightarrow\;\text{记录时刻表延续},
\]
其中第一箭头的反向蕴含、第二箭头的双向归约及时刻表的非周期性均在正文证明。记 $\vtwo{x}$ 为整数 $x$ 的 2-adic 赋值（$x$ 中 2 的幂次）。

\section{相位演算}\label{sec:phase}

\begin{definition}[事件与构型]\label{def:phase}
\textbf{事件}是带头处于状态 $D$ 且读 $0$ 的时刻。设磁带自某左侧哨兵之右为半无限，内容为 $1$-游程（极大 $1$-块）与 $0$-间隙交替，头位于某个 $0$ 上。四类构型（各类均可带尾部附加游程列）：
\begin{itemize}
\item $R(a,b)$：游程 $(a,b)$，间隙 $[1]$，头在右游程右侧的新鲜 $0$；
\item $S(x,y,z,\overline r)$：游程 $(x,y,z,\overline r)$，间隙 $[1,2,1,\dots]$，头在 $y$ 与 $z$ 间 2-间隙的左 $0$；
\item $G(a,b,\overline r)$：游程 $(a,b,\overline r)$，间隙 $[3,1,\dots]$，头在 $a$ 与 $b$ 间 3-间隙的中 $0$；
\item $L(r_1,r_2,r_3,\overline r)$：游程 $(r_1,r_2,r_3,\overline r)$，间隙 $[1,1,\dots]$，头在 $r_2,r_3$ 间的单 $0$。
\end{itemize}
\end{definition}

\begin{lemma}[相位规则]\label{lem:rules}
相邻两事件间，$M$ 恰执行下列转移之一（含精确步数）：
{\small
\[
\begin{array}{llll}
\text{(R)} & R(a,b)\to S(a{+}1,\,b{-}2,\,1) & 2b{+}3, & b\ge3;\\
\text{(S$\ge$)} & S(x,y,z,\overline r)\to S(x{+}1,y{-}2,z{+}1,\overline r) & 2y{+}3, & y\ge3;\\
\text{(S2)} & S(x,2,z,\overline r)\to G(x{+}1,\,z{+}1,\overline r) & 7;\\
\text{(S1)} & S(x,1,z,\overline r)\to L(1,\,x{+}z{+}3,\overline r) & 2x{+}z{+}14;\\
\text{(G$\ge$)} & G(a,b,\overline r)\to S(1,\,a{-}3,\,1,\,b{+}1,\overline r) & 2a{+}3, & a\ge4;\\
\text{(G3)} & G(3,b,\overline r)\to G(1,\,1,\,b{+}1,\overline r) & 9;\\
\text{(G2)} & G(2,b,\overline r)\to L(1,\,3,\,b{+}1,\overline r) & 16;\\
\text{(G1)} & G(1,b,\overline r)\to\HALT & 6;\\
\text{(L$\ge$)} & L(r_1,r_2,r_3,\overline r)\to L(r_1{+}1,r_2,r_3{-}1,\overline r) & 2r_2{+}5, & r_3\ge2;\\
\text{(L1)} & L(r_1,r_2,1,\overline r)\to S(r_1{+}1,r_2,\overline r) & 2r_2{+}5.
\end{array}
\]
}
(S1) 当尾 $\overline r$ 为空时结果记 $R(1,x{+}z{+}3)$；(L1) 当 $\overline r$ 空时结果为 $R(r_1{+}1,r_2)$。停机块 (G1) 以 $F$ 读 $0$ 结束；(G3) 是普通事件块，其落点 $G(1,\cdot)$ 在下一块停机（(G3)+(G1) 两块折叠共 15 步）。
\end{lemma}

规则由转移表逐步执行得到（(R) 的完整轨迹：$D$ 试探写入被 $E$--$F$--$E$--$A$ 回撤擦除，$A$ 左行填合间隙使首游程 $+1$，$B$ 清零一格生成新 2-间隙，$D$ 右行重建，净效果 $2b+3$ 步）。忠实性经两级机器验证：(i) 规则级 \texttt{rulecheck.mjs} 对参数网格（$a\le20$、$b\le10$、尾长 $0$--$2$）与裸机逐行对拍 $564/564$ 全等（含边界 $y\in\{1,2\}$、$a\in\{1,2,3\}$ 与停机情形 (G1)）；(ii) 轨迹级：相位引擎与裸机 oracle 在 956 个锚点（$n=2..957$）上 955 个逐字段全等（含 16 个停机锚点、停机时刻精确一致），唯一例外 $R(2,2)$（其后继 $b=2<3$ 在文法声明域之外，被引擎标记而非误算；直接模拟于 $+22$ 步停机）。早期版本携带折叠的 (G3)$\to\HALT$ 行与错误 (G2) 图像 $L(1,b{+}1,b{+}1)$，对拍该版本得 $531$ ok / $21$ bad（坏行恰为 (G2)/(G3)）；本文规则按裸机侧修正。转移表另经独立手工转写复核。

\begin{lemma}[停机判据]\label{lem:halt}
$M$ 停机当且仅当相位轨迹到达 $G(1,b,\overline r)$；等价地，当且仅当到达首头 $a\in\{1,3\}$ 的 $G$-态（(G3) 经一个事件块把 $G(3,\cdot)$ 送入 $G(1,\cdot)$）。
\end{lemma}
\begin{proof}
唯一停机转移为 $F$ 读 $0$。逐规则检查：十条规则的事件块中 $F$ 读到 $0$ 的情形仅在 (G1) 块内部出现；其余块中 $F$ 读 $1$ 或先到达下一事件。(G1)-态的唯一规则内来源是 (G3)。
\end{proof}

\begin{lemma}[扫掠算术]\label{lem:sweep}
(S$\ge$) 的连续应用可合并：$y$ 奇时 $j=\frac{y-1}{2}$ 步后 $y=1$、$x,z$ 各增 $j$；$y$ 偶时 $j=\frac{y-2}{2}$ 步后 $y=2$、$x,z$ 各增 $j$。同理 (L$\ge$) 合并为 $r_3\to1$。证明为对 $j$ 的直接归纳。
\end{lemma}

\section{锚点与序曲}\label{sec:prologue}

\begin{definition}[锚点]\label{def:anchor}
$P(n):=R(n,n)$，即 $0^\infty 1^n\,0\,1^n\,0^\infty$、头在右块右侧。空白带运行于第 54 步到达 $P(3)$（B 级：双模拟器逐步验证）。
\end{definition}

\begin{lemma}[减半复合]\label{lem:half}
设 $G(a,b,\overline r)$ 满足 $a$ 奇且 $a\ge5$。则
\[
G(a,b,\overline r)\;\longrightarrow\;G\!\left(\tfrac{a-1}{2},\,\tfrac{a-1}{2},\,b{+}1,\,\overline r\right),
\]
即头部对 $(a,b)$ 映为 $\bigl(\frac{a-1}{2},\frac{a-1}{2}\bigr)$，$b{+}1$ 作为新游程插入第三位，尾 $\overline r$ 逐字携带。
\end{lemma}
\begin{proof}
(G$\ge$)：$a\ge4$ 给 $S(1,\,a{-}3,\,1,\,b{+}1,\overline r)$。$a$ 奇 $\Rightarrow a{-}3$ 偶，由引理~\ref{lem:sweep}，$j=\frac{a-5}{2}$ 步扫掠后中游程到 $2$，两侧为 $1+j=\frac{a-3}{2}$ 与 $1+j=\frac{a-3}{2}$（$x=z=1$）。再用 (S2)：$G\bigl(\frac{a-3}{2}{+}1,\,\frac{a-3}{2}{+}1,\,b{+}1,\overline r\bigr)=G\bigl(\frac{a-1}{2},\frac{a-1}{2},b{+}1,\overline r\bigr)$。
\end{proof}

\begin{theorem}[序曲]\label{thm:prologue}
设 $n\equiv3\pmod 6$，$n\ge3$，$n{+}3=2^{m}w$，$w$ 奇，$m\ge1$。则从 $P(n)$ 出发的相位轨迹到达
\[
\Phi(w,m)\;=\;G\bigl(w{-}1,\;w{-}1,\;2w,\;4w,\;\dots,\;2^{m-1}w,\;2^{m}w{-}1\bigr),
\]
且途中头部遵循闭式 $a_i+1=2^{m-i}w$（$i=0,\dots,m$），无停机。
\end{theorem}
\begin{proof}
三段。
\emph{(i)} (R)：$R(n,n)\to S(n{+}1,\,n{-}2,\,1)$。$n$ 奇 $\Rightarrow y=n{-}2$ 奇，扫掠 $j=\frac{n-3}{2}$ 步后 $y=1$、$(x,z)=(n{+}1{+}j,\,1{+}j)$；用 (S1)：$R(1,\,x{+}z{+}3)=R(1,\,2n{+}2)$。
\emph{(ii)} (R) 于 $b=2n{+}2$：$S(2,\,2n,\,1)$；$2n$ 偶，扫掠 $j=n{-}1$ 步后 $y=2$、$(x,z)=(n{+}1,\,n)$；用 (S2)：$G(n{+}2,\,n{+}1)$。
\emph{(iii)} 令 $a_0=n{+}2=2^{m}w{-}1$（$a_0{+}1=2^m w$），$b_0=n{+}1$。由 $w$ 奇，$i<m$ 时 $a_i$ 奇且 $a_i\ge w{-}1\ge2$；对 $a_i\ge5$ 用引理~\ref{lem:half} 得 $a_{i+1}{+}1=\frac{a_i+1}{2}=2^{m-i-1}w$、$b_{i+1}=a_i$（旧首头成为第二游程）。$m$ 次后首头为 $w{-}1$，且第 $i$ 次复合插入的新游程为 $a_{i-1}{+}1=2^{m-i+1}w$，最后一次插入 $b_0{+}1=2^{m}w{-}1$。故游程列恰为
\[
\bigl(w{-}1,\;w{-}1,\;2w,\;4w,\;\dots,\;2^{m-1}w,\;2^{m}w{-}1\bigr)=\Phi(w,m). \qedhere
\]
\end{proof}

\begin{remark}
$a_i$ 奇（$i<m$）因 $a_i{+}1=2^{m-i}w$ 含偶因子；$a_m=w{-}1$ 偶，故级联恰在 $m$ 步后自然终止于 $\Phi$，无 (G1)/(G3) 风险（停机族首头 $\in\{1,3\}$：$w{-}1=1$ 需 $w{=}2$、$w{-}1=3$ 需 $w{=}4$，均与 $w$ 奇矛盾）。定理为纯代数，另经 $n\le6003$ 全体（1001 例）与 $m\le1200$ 样本机器复核（B 级）。
\end{remark}

\section{模 3 台账与 2-adic 出生}\label{sec:ledger}

本节给出两个\textbf{不依赖假设 M} 的定理——全轨道模 3 台账（定理~\ref{thm:mod3}）与停机等价（定理~\ref{thm:haltpow}）——及出生守卫污染的见证级命题（命题~\ref{prop:depth1}、\ref{prop:cascade}）。二者合起来：模 3 侧被无条件封死，停机排除的全部剩余内容浓缩为一句关于出生值的 2-adic 陈述。

\begin{definition}[台账不变量 $I_0$]\label{def:I0}
对构型 $C$，令 $\Sigma(C)$ 为其游程总和。$I_0$ 为满足以下条件的构型集：(i) $\Sigma\equiv0\pmod3$；(ii) 域：$R(a,b)$ 需 $b\ge3$，$S(x,y,z,\overline r)$ 需 $y\ge1$，$G(a,b,\overline r)$ 需 $a\ge2$，$L(r_1,r_2,r_3,\overline r)$ 需 $r_2\ge3$ 且 $r_3\ge1$；(iii) 边界守卫 BG1：$G(2,b,\overline r)$ 需 $b\ge2$ 且 $b\equiv2\pmod3$；BG2：$S(1,2,z,\overline r)$ 需 $z\equiv1\pmod3$。
两条守卫均为封闭\textbf{被迫}：BG2 经 (S2) 生出 $G(2,z{+}1)$，子须过 BG1；BG1 经 (G2) 的引擎式 $\Sigma$-增量 $\Delta=b{+}1$ 须 $\equiv0$（裸机层该边恒 $\Delta=+3$，见注~\ref{rem:g2}）。联合 $\Sigma\wedge\mathrm{BG2}$ 下三游程的 $S(1,2,z)$ 整类为空（$z\equiv0$ 与 $z\equiv1$ 矛盾）；$k\ge4$ 的 $S(1,2,z,\overline r)$ 因尾自由仍非空，而轨道恰走后者（(G$\ge$) 于 $a=5$ 的子 $S(1,2,1,b{+}1,\overline r)$，$z=1$）。
\end{definition}

\begin{lemma}[$G(2,\cdot)$ 的唯一可达值]\label{lem:g22}
文法可达的 $G(2,b,\overline r)$ 恰为 $b=2$ 者（含几何尾族 $\Phi(3,m)$）。
\end{lemma}
\begin{proof}
$G$-态的唯一规则内来源是 (S2)：$S(x,2,z,\overline r)\to G(x{+}1,z{+}1,\overline r)$，首头为 $2$ 迫使扫掠终态 $x=1$。而 $x=1$ 的 $S$-入口只能来自 (G$\ge$)（(R) 生成 $x=a{+}1\ge2$，(L1) 生成 $x=r_1{+}1\ge2$，(S$\ge$) 扫掠中 $x$ 单调增），即 $S(1,a{-}3,1,b{+}1,\overline r)$，其中 $z=1$ 固定；扫掠至 $y=2$ 后仍有 $x=1$ 者仅零步扫掠（$a=5$）一例，得 $S(1,2,1,b{+}1)\to G(2,2,b{+}1,\overline r)$。
\end{proof}

\begin{theorem}[模 3 台账；无条件于假设 M]\label{thm:mod3}
\textbf{(a) 封闭性}：$I_0$ 的每个成员经引理~\ref{lem:rules} 十条规则之任一单步后仍属 $I_0$，唯一例外是停机形态 $G(3,b,\overline r)$（被分类，非违例）。\textbf{(b) 种子}：空白带轨道于第 54 步进入 $R(3,3)\in I_0$（$\Sigma=6$）。\textbf{(c) 结论}：自第 54 步起 $M$ 的一切相位（文法）构型满足 $\Sigma\equiv0\pmod3$。从而：每个重入 $R(1,b)$ 有 $b\equiv2\pmod3$；每个锚 $R(n,n)$ 有 $n\equiv0\pmod3$；任何停机只能发生在 $\Sigma\equiv0$ 的 $G(3,b,\overline r)$；一切 $\Sigma\not\equiv0$ 的构型——包括 Conway 型迫使链入口 $S(3,1,3)$（$\Sigma=7$）——构成整半空间不可达。$I_0$ 不排除停机（$G(3,\cdot)$ 停机形态可在 $I_0$ 内）；停机排除的全部剩余内容在 2-adic 出生侧（定理~\ref{thm:haltpow}）。
\end{theorem}
\footnote{停机段粒度：折叠块 (G3)$+$(G1)（共 15 步）于 $+9$ 步处含中间事件 $G(1,1,b{+}1,\overline r)$，即 (G3) 的目标——在 $I_0$ 的 $G$-域（$a\ge2$）之外，但其 $\Sigma$ 仍 $\equiv0$（(G3) 的 $\Sigma$-恒等式 $\Delta=0$；裸机抽验一致）。定理~\ref{thm:mod3}(c) 的全称陈述因此钉在\textbf{相位（文法）构型}上。}

证明地位（A$^-$，封闭部分）：封闭 = 逐边初等仿射恒等式 + 守卫转移——(R)：$(a{+}1)+(b{-}2)+1=a+b$；(S$\ge$) 与 (L$\ge$)：平移；(S2)：$(x{+}1)+(z{+}1)=x{+}z{+}2$；(S1) 两支各 $+3$；(G$\ge$)：$1+(a{-}3)+1+(b{+}1)=a+b$；(G3)：$\Delta=0$；(G2)：见注~\ref{rem:g2}；(L1) 两支各 $+1$——逐边可手工复核，另经机器全查复核：作者侧边界穷举 219,398 个 $I_0$ 成员步 $+$ $2^{40}$ 量级随机成员步（断言总数随随机抽样摆动约 $\pm60$，确定性核心逐位可复现）；独立侧异构封闭网格 20,508 个成员步，游程数至 $k\le10$，含 $a=0$ 探针。封闭性对游程数的归约由局部性承载（每条规则只改写前三个游程，两条守卫只读前三个游程）。裸机逐事件忠实性经独立复核（十条规则，边界网格 $+$ $2^{24}$ 量级随机）；种子三源（双模拟器逐步 $+$ 独立裸机解析）。轨道侧为冗余确认（归纳已给全时域结论）：$10^7$ 相位事件零违例（406 锚全 $n\equiv0$、1,485 重入全 $b\equiv2$）；独立复核 $5\times10^5$ 事件零违例（91 锚、332 重入）；轨迹库 $1{,}357{,}014$ 态零违例（$R$-族分解 $825=547{+}150{+}128$，全部 $a{+}b\equiv0$）。文内的 112 个 $G(3,\cdot)$ 停机形态为封闭扫描网格的枚举计数，非轨道实测（轨道上 $G(3,\cdot)$ 计数为 $0$）。

\begin{remark}[(G2) 的裸机逼真度域]\label{rem:g2}
裸机真值：(G2) 边的真实映射为 $G(2,b,\overline r)\to L(1,3,b{+}1,\overline r)$，对一切 $b\ge2$ 成立（$b\in\{2,3,5,8,11,14\}$、三种尾裸机验证；与本文规则表逐字一致，rulecheck 564/564 亦含各 $b$ 的 (G2) 行），其 $\Sigma$-增量恒为 $+3$——故模 3 台账在裸机层于 $I_0$ 全域保持，\textbf{不依赖 BG1}（BG1 在裸机恒等式下反成多余，仅引擎式记账需要）。台账分析套件的推导引擎携带简化图像 $L(1,b{+}1,b{+}1,\overline r)$，其与裸机一致当且仅当 $b=2$；由引理~\ref{lem:g22}，$b=2$ 是唯一文法可达值，故失真位点 $G(2,b\neq2)$ 不可达，且即使落于其上（若自文外构型启动）$\Delta=+3$ 亦保持台账。据此，``十条规则 $=$ 裸机''的读法须限定为：规则表在文法域上逐事件裸机忠实；文法外层（$G(2,b\neq2)$）上引擎为简化式，台账不受影响。
\end{remark}

\begin{definition}[S-偶出生]\label{def:birth}
到达构型 $S(x,y,z,\overline r)$ 且 $y$ 为偶数（$y\ge2$）的事件称为一次 \textbf{S-偶出生}（该偶扫掠的入口），其\textbf{出生值}为
\[
A \;=\; x+\tfrac{y}{2}+1 .
\]
\end{definition}

\begin{theorem}[停机等价：2-幂出生]\label{thm:haltpow}
$M$（自空白带）停机，当且仅当其相位轨道上某次 S-偶出生值 $A=x+y/2+1$ 为 2 的幂；等价地 $A\in\{4,8,16,\dots\}$（恒有 $A\ge3$，因 $x\ge1$、$y\ge2$；$A=3$ 的扫掠终态首头为 $2$，经 (G2) 入非停机通道）。
\end{theorem}
\begin{proof}
先注意\textbf{扫掠不变量}：(S$\ge$) 保持 $A$，因 $(x{+}1)+\frac{y-2}{2}+1=x+\frac y2+1$；出生扫掠终止于 $S(x',2,z')$（$x'=x+\frac{y-2}{2}$），经 (S2) 落入首头 $x'{+}1=A-1$ 的 $G$-态。
（$\Longleftarrow$）设某次出生 $A=2^t$（$t\ge2$）。其扫掠终态 $G$-态首头为 $2^t-1$（奇，$\ge3$）。$t=2$ 时即停机形态 $G(3,\cdot)$。$t\ge3$ 时首头 $\ge7$ 且奇，由引理~\ref{lem:half} 逐次减半复合：首头链 $2^t-1\to2^{t-1}-1\to\cdots\to7\to3$（尾游程逐字携带；链上首头皆形如 $2^s-1\ge3$，永不取 $5$ 等其他奇值，亦不落 $G(2,\cdot)$ 通道），终于停机形态 $G(3,\cdot)$，由引理~\ref{lem:halt} 停机。
（$\Longrightarrow$）设 $M$ 停机。由引理~\ref{lem:halt}，轨迹到达首头 $3$ 的 $G$-态；$G(3,z{+}1,\overline r)$ 的唯一规则内来源是 (S2) 作用于 $S(2,2,z,\overline r)$。该 $S$-态 $y=2$，位于某次偶扫掠之内（或自身即入口）；扫掠不变量给出该扫掠入口处的出生值 $A=2+\frac22+1=4=2^2$，为 2 的幂。
\end{proof}

证明地位（A$^-$，双向）：两个方向均为初等代数，立足于规则忠实性与停机形态唯一性（首头 $\in\{1,3\}$ 的 $G$-态只有 (G3)/(G1) 两条出路，经裸机广检）；规则忠实性为 B 级对拍（rulecheck 564/564、956 锚点）加独立裸机逐事件复核。轨道实证（B 级）：前 $10^7$ 相位事件共 2,908 次 S-偶出生，0 次为 2 的幂；独立复核 $5\times10^5$ 事件 647 次出生，0 次。假设 M 的全部内容浓缩为``永无 2-幂出生''这一句话。

\begin{proposition}[深度 1 出生守卫不封闭]\label{prop:depth1}
令 $I_1$ 为 $I_0$ 加上逐出生点的同形代数守卫（出生值奇部 $\ge3$；$G(2,\cdot)$ 的 $f_2$ 守卫；$A=3$ 位点的 $w\equiv2\pmod3$；精确式见工件 \texttt{invariant-casecheck} 系列）。则 $I_1$ 含全部真实轨迹（轨迹库 $1{,}357{,}014$ 态零违例），但 $I_1$ \textbf{不封闭}：声明域（$S/L$ 尾长 $\le6$）内的封闭搜索得 1,255 个违例，全部属同一类——奇 $r_2$ 的 $L$ 经 (L1) 武装其 S-奇子，而子的局部守卫对其下一步出生盲。
\end{proposition}

\begin{proposition}[出生守卫的逐深度污染；见证级]\label{prop:cascade}
对每个已测深度 $d\in\{1,2,3\}$，存在 $I_0$ 内构型，其管线前 $d-1$ 次 S-偶出生非 2 的幂、第 $d$ 次出生为 2 的幂，且整链真停机：
\begin{center}
\begin{tabular}{clllr}
\toprule
深度 & 见证构型 & 出生序列 & 相位事件 & 裸机步数\\
\midrule
$d{=}1$ & $L[1,3,1,3,1]$ & $8,4$ & 13 & 195\\
$d{=}1'$ & $S[1,1,1,1,1,1]$ & $8,4$ & 15 & 230\\
$d{=}1''$ & $G[4,2,7,5]$ & $16,8,4$ & 33 & 728\\
$d{=}2$ & $S[3,2,25,43,23,12,4,2]$ & $5,64$ & 152 & 10,313\\
$d{=}3$ & $S[12,28,23,42,14,22,38,31]$ & $27,71,128$ & 542 & 71,856\\
\bottomrule
\end{tabular}
\end{center}
五见证均经独立裸机复现（逐事件磁带与相位构型完全一致、裸机停机确认、全链 $\in I_0$；另：步数上限伪影见证 $S[39,34,5,30,43,19,37]$ 于大上限下 442 事件停机，非反例）。且\textbf{剩余类细化不可能}：坏出生 $\{2^t\}\bmod12\in\{4,8\}$，而轨道自身的 S-偶头恰占据同类（$A\bmod12$ 直方图覆盖 $\{0,\dots,9\}$，其中 $A\equiv4$ 共 9,282 例、$A\equiv8$ 共 396 例），类内切分恰为 $\nu_2$ 之别——不可约 2-adic。坏出生的进入边为\textbf{出生进入边族}：(L1) 为主（$I_1$ 的 1,255 个违例与 $d=2$ 见证均经之进入），(R) 亦出现（$d=3$ 见证经 (R)$\to S_3$ 进入）；污染机制为 (S1)/(L1) 管线的逐层平移（注~\ref{rem:translate}）。
\end{proposition}

\begin{remark}[一般深度 $d\ge4$：结构平移机制，未证明]\label{rem:translate}
逐层修复（深度 $d$ 守卫 := ``管线前 $d$ 次出生均非 2-幂''）后，同一形状的 (S1)/(L1) 边在下一个 $L$-轮次复现一层更深：每修一层，同形管线平移一层。据此\textbf{预期}对每 $d\ge4$ 亦有见证（随机搜索于各深度稳定产出新见证），但逐 $d$ 见证的存在性恰是待证命题，本文不主张；``任意有限守卫深度不可封闭''以见证级（命题~\ref{prop:cascade} 与本机制）表述，非定理。
\end{remark}

\begin{remark}[与猜想~\ref{thm:noprec} 的关系]\label{rem:noprec2}
命题~\ref{prop:cascade} 不证明猜想~\ref{thm:noprec}，亦不改变其（猜想）地位；它给出的是\textbf{猜想之外的强化方向}——模板语言扩至``剩余类 $+$ 代数头形式守卫''后，有限守卫深度的封闭仍逐深度失败——的见证级支持（B 级）。迫使边不在模 3 侧（定理~\ref{thm:mod3} 已把该侧完全封死），而在 2-adic 出生位点：进入边为出生进入边族（(L1) 主、(R) 亦见），机制为 (S1)/(L1) 管线平移（注~\ref{rem:translate}）。
\end{remark}

\section{地标、类与轮定律}

固定奇单位 $w\ge3$。

\begin{definition}[地标与类]\label{def:landmark}
对 $\alpha\equiv1\pmod6$，$\alpha\ge7$，$t\ge9$，松弛 $s\ge0$，令 $E=2^{t-3}w$：
\[
L(\alpha,t;s)=G\bigl(\alpha{+}3E,\;\alpha{-}1{+}5E,\;2^{t}w,\;2^{t+1}w,\;\dots,\;2^{t+s}w,\;2^{t+s+1}w{-}1\bigr),
\]
其\textbf{类}为 $h=\vtwo{\alpha+1}$（$\alpha{+}1=2^{h}\omega$，$\omega$ 奇）。
\end{definition}

\begin{lemma}[类参数化]\label{lem:param}
固定 $h\ge1$：满足定义~\ref{def:landmark} 的 $\alpha$ 恰为 $\alpha=2^{h+1}(3c+\rho_h)+2^{h}-1$（$c\ge0$），其中 $\rho_h\in\{0,1,2\}$ 为使 $\alpha\equiv1\pmod3$ 的唯一剩余。
\end{lemma}
\begin{proof}
$\alpha{+}1=2^h\omega$、$\omega$ 奇 $\Rightarrow\omega=2(3c{+}\rho){+}1$；三个 $\rho$ 给出 $\alpha\bmod3$ 的三个剩余，恰一者合于 $\alpha\equiv1$。
\end{proof}

\subsection*{轮定律与强归纳架构}

\begin{assumption}[主控封闭，M]\label{ass:M}
存在函数 $\Delta\alpha,\Delta t:\mathbb N_{\ge1}\to\mathbb N$，使对\textbf{每个}类 $h\ge1$ 与每个类 $h$ 地标 $L(\alpha,t;s)$（任意奇单位 $w\ge3$、任意外层参数 $c,u$，$t\ge2h{+}16$）：
(i) $s\ge\Delta t(h)$ 时轨迹到达下一地标 $L(\alpha{+}\Delta\alpha(h),\,t{+}\Delta t(h);\,s{-}\Delta t(h))$，且 $\Delta\alpha(h)\equiv0\pmod6$、$\Delta\alpha(h)\le 2h{+}6$、$\Delta t(h)\ge2$；
(ii) $s<\Delta t(h)$ 时到达锚点出口 $n'=2^{t+s+1}w+c'$，$c'\equiv3\pmod6$，$c'\ge3$。
两条路径均不经过 $\HALT$；推导（自阶梯字面量 $(A_{j_0(h)},\tau_{j_0(h)},h)$ 出发的符号推导）为单路径并於 $O(\mathrm{poly}(\log h))$ 事件内终止，且其形式集\textbf{$h$-均匀}：仅经阶梯字面量与修正尾的指数偏移依赖 $h$（与 en 版论文的主控封闭假设同一形式化）。
\end{assumption}

M 的\textbf{内容}由主控段承担：自包含符号推导（引擎 \texttt{sym5}，复杂度 $O(\mathrm{poly}(\log h))$，分支单路径）。其已验范围（B 级）：$h\le2000$ 穷尽、抽样至 $1.2\times10^7$；幅度类分支由支配引理无条件覆盖（下节）。

\begin{lemma}[级联]\label{lem:cascade}
类 $h$ 地标 $L(\alpha,t;s)$（$t\ge h{+}4$，$\omega=(\alpha{+}1)/2^{h}$）的前 $h$ 个 $G$-事件为减半复合，第 $h$ 个落点为 $G(\Psi{-}1,\Psi{-}1,2\Psi,\dots,2^{h-1}\Psi,\,\beta,\,2^{t}w,\dots)$，其中 $\Psi=\omega+3\cdot2^{t-3-h}w$（奇），$\beta=\alpha{+}5E$。
\end{lemma}
\begin{proof}
首头 $a_0=\alpha{+}3E$：$a_0{+}1=2^{h}\omega+3\cdot2^{t-3}w=2^{h}\Psi$，故奇；由引理~\ref{lem:half} 逐次减半得 $a_i{+}1=2^{h-i}\Psi$，$i<h$ 时奇。插入游程依次为 $a_{i-1}{+}1=2^{h-i+1}\Psi$，形成 $\Psi$-几何尾；第 $h$ 次复合时原第二游程 $\alpha{-}1{+}5E$ 按引理~\ref{lem:half} 的 $b{\mapsto}b{+}1$ 插入规则成为新第三游程 $\beta=\alpha{+}5E$（逐游程等式经 \texttt{entry.mjs} 于 $h\le60$ 全体验证，B 级）。
\end{proof}

\begin{lemma}[阶梯与良基性]\label{lem:ladder}
定义梯 $(A_0,\tau_0)=(7,9)$，$(A_{j+1},\tau_{j+1})=(A_j{+}\Delta\alpha(H_j),\,\tau_j{+}\Delta t(H_j))$，$H_j=\vtwo{A_j+1}$。则：
(i) $A_j\le\tau_j^2$（$j\ge8$）时 $A_{j+1}\le\tau_{j+1}^2$ 持续成立；
(ii) 凡 $\tau_j\le h{-}1$ 的位置有 $H_j\le2\log_2 h<h$（$h\ge16$）。
\end{lemma}
\begin{proof}
(i) 设 $A_j\le\tau_j^2$。由 M 的带界 $\Delta\alpha\le2H_j{+}6$ 与 $H_j\le\log_2(A_j{+}1)\le2\log_2\tau_j$：
$A_{j+1}\le\tau_j^2+4\log_2\tau_j+6\le\tau_j^2+4\tau_j+4=(\tau_j{+}2)^2\le\tau_{j+1}^2$（末步用 $\Delta t\ge2$ 与 $4\log_2\tau{+}2\le4\tau$，$\tau\ge2$）。
(ii) $H_j=\vtwo{A_j{+}1}\le\log_2(A_j{+}1)\le2\log_2\tau_j\le2\log_2(h{-}1)<h$。
基例：(i) 于 $j\in\{8,\dots,200000\}$ 与 (ii) 于所涉区间的前 $200{,}000$ 个位置（覆盖一切 $\tau_j\le10^6$）经穷举验证（B 级；与 en 版论文同一枚举）。$j<8$ 的位置（$\tau_j\le23$）直接包含于该枚举。
\end{proof}

\begin{theorem}[轮定律：条件证明]\label{thm:round}
假设 M 成立，则对一切类 $h\ge1$、一切类 $h$ 地标，轮定律（M 的 (i)/(ii) 两分支）成立，且全程无停机。
\end{theorem}
\begin{proof}
对 $h$ 强归纳。$h\le15$：直推表逐类单路径符号推导（B 级，1024 行表之子集）。$h\ge16$：类 $h$ 地标的一轮 = 级联（引理~\ref{lem:cascade}，代数）+ 入口（至 $\Psi$-尺度 $L(7,9)$ 型，读窗指数 $\le10<h{-}1$，修正尾逐字携带；\texttt{entry2.mjs} $m$-无关性 1156/1156，B 级）+ 阶梯前缀（引理~\ref{lem:ladder}(ii)：所历位置类 $<h$，由归纳假设其轮定律成立）+ 主控段（假设 M 于类 $h$）。合成的输出即 M 的 (i) 或 (ii)，无 (G1)/(G3) 事件。带界 $\Delta\alpha\le2h+6$ 已含于 M，随归纳自洽。
\end{proof}

\section{主控段与签名动力学：M 的归约}\label{sec:reduce}

主控段推导中每个分支判定为两类之一：\textbf{幅度比较}（游程大小与阈值）与\textbf{奇偶判定}。

\begin{lemma}[支配，D]\label{lem:dom}
每个幅度比较的裕度为 $M=W+Q-\theta$，其中 $w$-项 $W\ge3\cdot2^{2h-52-2\log_2 h}$，余项 $|Q|<2^{h+5}$，$\theta\le4$。故 $h\ge82$ 时 $M>0$ 且优势每单位 $h$ 翻倍；$h\le81$ 逐点验证。
\end{lemma}
\begin{proof}[证明概要]
$W$ 的下界来自读窗与类界引理（比较对象中至少一方含 $w$-尺度几何项，指数 $\ge-65-3-2\log_2 h$）；$|Q|$ 的上界来自 $d$-侧位数的仪表记录（$\le h{+}4$，12000 条记录无一违反——\textbf{此上界为 B 级（仪表穷举），非纸面证明}，故本引理整体降为 A$^-$）。$2h-52-2\log_2h>h+5\iff h>57+2\log_2h$，$h\ge82$ 成立。
\end{proof}

\begin{lemma}[剖面函数性]\label{lem:profile}
主控段的完整事件序列（因而结果 $\Delta\alpha(h),\Delta t(h)$ 或出口参数）是阶梯窗口签名的函数：签名相同的两次主控推导事件序列逐事件相同。剖面 = 级联头值的 2-adic 展开（有限下降，于类深度处终止）。
\end{lemma}
\begin{proof}[证明概要]
由 D，幅度判定在全域被支配（恒取正支），无自由度；奇偶判定由窗口数据决定（12{,}000 条仪表记录零定义域分裂）。级联头的 2-adic 展开有限，故常数剖面（需无穷下降）不可实现。
\end{proof}

\begin{definition}[签名动力学]\label{def:sig}
固定 K-切片模型：状态为窗口签名 $=$（最近 $\le40$ 个梯位置的 $(A_j\bmod 2^{K},H_j)$，预算 $b=h{-}1{-}\tau_{j_0}$——预算无界，故完整状态空间非有限；\textbf{残差核}（固定预算、只留签名）为有限状态系统）。演化一阶：下一签名由当前签名决定；唯一的非马尔可夫断点是 $\vtwo{}$-截断（真深度 $\ge K$），即\textbf{记录事件}。模型忠实性属 B 级（\texttt{sigdyn.mjs}/\texttt{brent32.mjs} 与符号引擎对拍）。
\end{definition}

\begin{theorem}[归约定理]\label{thm:reduce}
在 K-切片模型内：假设 M（对一切类）成立当且仅当自主游走的记录时刻表无限延续（每个深度级的主控推导在其被到达时闭合）。
\end{theorem}
\begin{proof}
($\Leftarrow$) 各深度级的主控推导只依赖窗口签名与预算（引理~\ref{lem:profile}）；时刻表延续即每个被到达的类都有已闭合的推导，故 M 逐类成立。
($\Rightarrow$) 若时刻表在深度 $d$ 终止（某级不再到达），则该级之后无新类，M 仅需已有类。到达类集合由 $(A_j,\tau_j)$ 递推决定；时刻表终止将迫使 K-切片残差核（$K'=\max(d,K)$ 位）进入环，环内深度被压于 $K'$ 以下永续——与时刻表\textbf{在已验切片上的}爬升矛盾。此方向对未观测深度的外推为\textbf{条件性}（有限观测不能排除未来抑制；其无条件形式即定理~\ref{thm:nocycle} 所述 K$>$32 猜想）。两种情形都把 M 的内容完全转移到时刻表上。
\end{proof}

\begin{theorem}[周期排除（32-位切片证书；K$>$32 为猜想）]\label{thm:nocycle}
签名动力学的 32-位残差核无环且终止：Brent 环检测于 $1.25\times10^9$ 步无环（核终止于坏残差 $\alpha\equiv2^{32}{-}1$，即记录-32 事件）；$2.4\times10^6$ 个梯残差两两互异；记录时刻表在已验域内每级 $\times4$ 爬升（实测 $17@19024\to19@76096\to21@304383\to22@1217531\to24@4870122\to\ge32@1.25\times10^9$）。故 $K\le32$ 内不存在有限闭环。任何 $K$-位环将把记录永久压于深度 $K$ 以下；``全深度非最终周期''（$K>32$）为\textbf{猜想}——爬升律在已验切片外是测量而非证明。
\end{theorem}
（证据为 B 级计算；``环不存在则 M 不可有限化''的转移为代数。）此即：M 的剩余内容 $=$ 游走的 $\vtwo{\cdot}$-混沌，本架构不可约的硬核。

\begin{definition}[粗模板谓词]\label{def:template}
\textbf{粗模板}谓词 $I$：构型的隶属仅依赖游程序列的（a）固定模数 $M_0$ 的逐游程剩余类 $r_i\bmod M_0$、（b）有限组线性不等式 $\lambda_1 r_1+\dots+\lambda_k r_k\ge\theta$、（c）游程数 $k$ 属于有限集。（即：对 $M_0$ 逐游程同余且满足同组不等式的构型不加区分。）
\end{definition}

\begin{conjecture}[必要精度：粗不变量架构的不可能性]\label{thm:noprec}
不存在粗模板谓词 $I$ 同时满足：(i) $P(3)\in I$；(ii) $I$ 对引理~\ref{lem:rules} 的十条规则封闭；(iii) $I$ 排除 $G(1,\cdot)$ 与 $G(3,\cdot)$。
\end{conjecture}
\begin{remark}[证据与已核部分]\label{rem:noprec}
(a) \textbf{反例链（A 级，规则级机器复核）}：停机族经少数几步可达自 $R(1,9)$ 与 $R(1,12)$ 型——
\[
S(3,1,3)\xrightarrow{\text{(S1)}}R(1,9)\xrightarrow{\text{(R)}}S(2,7,1)\xrightarrow{\text{扫掠}\times3}S(5,1,4)\xrightarrow{\text{(S1)}}R(1,12)\xrightarrow{\text{(R)}}S(2,10,1)\xrightarrow{\text{扫掠}\times4}S(6,2,5)\xrightarrow{\text{(S2)}}G(7,6)\xrightarrow{\text{(G$\ge$)}}S(1,4,1,7)\xrightarrow{\text{(S$\ge$)}}S(2,2,2,7)\xrightarrow{\text{(S2)}}G(3,3,7)\xrightarrow{\text{(G3)}}G(1,1,8).
\]
（本文早期版本写出的 $R(3,9)/S(4,7,1)$ 链有误，已废弃。）故任何封闭谓词必须在 $R(1,\cdot)$ 族内把 $9/12$ 型与真实轨道的重入型分开。
(b) \textbf{模 3 护城河（已定理化）}：定理~\ref{thm:mod3}(c) 无条件给出真实轨道自第 54 步起一切相位构型 $\Sigma\equiv0\pmod3$（原 15 万事件 16,826 个重入零违例的 B 级实测降为冗余确认；锚点序曲的 $b=2n{+}2$、$n\equiv3\bmod6$ 给 $b\equiv8\equiv2\bmod3$ 与之一致）——而 (a) 的反例链自 $S(3,1,3)$（$\Sigma=7$）起\textbf{整链} $\Sigma\equiv1\pmod3$（(S1) $+3$、其余边 $\Delta\equiv0$，同余类沿链保持），入口起即被定理~\ref{thm:mod3} 无条件排除。但 $G$-首头的轨道分布覆盖全部模 3 类（实测直方图 $2{:}30348$、$0{:}1408$、$1{:}1330$），故该条件不能经其余形式封闭——完整的不可能性证明须在全体形式/规则上组合此类模约束，本轮未完成（猜想维持原级；强化方向见注~\ref{rem:noprec2}）。
(c) 不变量合成实验（\texttt{invsyn.mjs}，否定结果）：未找到任何粗模板不变量。深层原因（启发式）：真实轨迹上 $G$-头在 2-adic 意义下向停机族累积（$\nu_2(a{-}1),\nu_2(a{-}3)$ 达 9 已现于轨迹），固定模数同余对 $\vtwo{a-1}/\vtwo{a-3}$ 盲。
\end{remark}

\begin{corollary}[双架构收敛（条件于猜想~\ref{thm:noprec}）]\label{cor:conv}
条件于猜想~\ref{thm:noprec} 与定理~\ref{thm:nocycle} 的 32-位切片证书：任何不重算假设 M 内容的封闭谓词都不足以排除停机；任何足以排除者必携带 $\vtwo{\cdot}$-混沌的全部信息。故在此两类架构内，假设 M 不可消除。
\end{corollary}

\section{装配与主定理}

\begin{theorem}[主定理，条件形式]\label{thm:main}
假设 M 成立，则 $M$ 从空白带启动永不停机。
\end{theorem}
\begin{proof}
空白带第 54 步到 $P(3)$（定义~\ref{def:anchor}）。对锚点下标 $k$ 归纳，维持不变量 $6\mid n_k{+}3$：$n_0=3$ 成立；设 $6\mid n_k{+}3=2^{m_k}w_k$，则 $3\mid w_k$。由定理~\ref{thm:prologue} 轨迹入 $\Phi(w_k,m_k)$；由定理~\ref{thm:round}（级联、阶梯、主控）诸轮后经 (ii) 出口到锚点 $n_{k+1}=2^{m_k}w_k+c'$，$c'{\equiv}3\pmod6$、$c'{\ge}3$。于是 $n_{k+1}{+}3=2^{m_k}w_k+(c'{+}3)$，两加项均可被 6 整除（$2\mid2^{m_k}$、$3\mid w_k$；$6\mid c'{+}3$），不变量保持且 $n_{k+1}\ge n_k{+}6$。全程无 $\HALT$（定理~\ref{thm:round} 与停机判据引理~\ref{lem:halt}）。锚点列严格递增无上界，运行永续。
\end{proof}

\begin{corollary}[无条件包络（验证域；经验成分显式标注）]\label{cor:env}
不假设 M：对一切类 $\le2000$（穷尽验证；抽样一致至 $1.2\times10^7$）的运行，$M$ 不停机。步数换算为双重指数视界时含\textbf{经验成分}：记录时刻表的 $\times4$/级爬升为实测律（定理~\ref{thm:nocycle}），非无条件定理；类 $h$ 首达时间 $\ge2^{c\cdot2^{h}}$ 的纸面下界（由 $A_j\le\tau_j^2$ 与带界 $\Delta\alpha\le2h{+}6$ 组装）列为后续。另，停机分支在代数族上无条件不可达：轨道序曲/级联/地标族的一切 $G$-头部满足 $a{+}1=2^{j}\cdot(\text{奇}\ge3)$，故 $a\notin\{1,3\}$（$a{+}1\in\{2,4\}$ 的奇部为 $1<3$）；由引理~\ref{lem:dom} 该结论对 $\pm1$ 扰动稳健（2980/2980 反事实闭合安全；对抗剖面被 2-adic 展开有限性排除且从不停机）。中间态/主控态/残形态的排除由 2-adic 审计承载（273,654 形式零危险），其对全类的推广条件于假设 M 的 $h$-均匀条款。另由定理~\ref{thm:mod3} 与 \ref{thm:haltpow}，无条件包络的形状如下精化：一切 $\Sigma\not\equiv0\pmod3$ 的构型自第 54 步起不可达（整半空间，不依赖任何已验类域），且``$M$ 永不停机''等价于``永无 2-幂 S-偶出生''——后一等价式的已验域为前 $10^7$ 相位事件（2,908 次出生零 2-幂，B 级）。
\end{corollary}

\section{台账}

\begin{center}
\begin{tabular}{clp{8.2cm}}
\toprule
级 & 含义 & 本文条目\\
\midrule
A & 纸面代数 & 引理~\ref{lem:sweep}、\ref{lem:half}、定理~\ref{thm:prologue}、引理~\ref{lem:param}、\ref{lem:cascade}（尾簿记部分 B）、\ref{lem:ladder}（基例 B）、定理~\ref{thm:round}（给定 M）、\ref{thm:reduce}（模型内，$\Rightarrow$ 方向含外推条款）、\ref{thm:nocycle}（K$\le$32 部分与转移部分）、定理~\ref{thm:main}\\
A$^-$ & 纸面骨架 + 仪表穷举上界 & 引理~\ref{lem:dom}（$|Q|<2^{h+5}$ 来自 12,000 条仪表记录）；定理~\ref{thm:mod3}（封闭 = 仿射恒等式 + 守卫转移，机器全查复核 + 裸机逐事件忠实性独立复核）；定理~\ref{thm:haltpow}（双向代数，立足于规则忠实性与停机形态唯一性）\\
B & 声明域穷尽 & 规则级对拍（\texttt{rulecheck} 564/564）与轨迹忠实性（956 锚点：955 一致 + $R(2,2)$ 域外，16 停机锚精确）；序曲复核（1052）；级联/入口（327；\texttt{entry2} 1156/1156）；轮律表（1024+976）；主控扩展（至 $1.2\times10^7$ 抽样）；残局格（5890）；2-adic 审计（273{,}654 形式）；反事实（2980）；Brent（$1.25\times10^9$ 步）；时刻表数据；护城河实测（$R(1,b)\equiv2\bmod3$、$G$-头模 3 直方图）；台账全查（作者 219,398 成员步 + 独立异构 20,508 步 $k\le10$；$10^7$ 事件零违例、2,908 出生零 2-幂；独立 $5\times10^5$ 事件 647 出生零幂；库 $1{,}357{,}014$ 态零违例）；级联见证（13/15/33/152/542 事件裸机复现）与 $I_1$ 封闭搜索（1,255 违例一类）、$A\bmod12$ 直方图（9,282+396）——命题~\ref{prop:depth1}、\ref{prop:cascade} 为\textbf{见证级}\\
C & 依赖假设/猜想 & 假设 M（$\equiv$ 记录时刻表延续，定理~\ref{thm:reduce}）；猜想~\ref{thm:noprec}（粗模板不可能性）；K$>$32 非周期性\\
\bottomrule
\end{tabular}
\end{center}

\end{document}
